\subsection{Substitution Ciphers}


\content{
        Last time we talked about shift ciphers, and determined that they are
        not very secure (there are only 26 possible secret keys). What if
        instead we choose some random mapping, like the following

        \texttt{ABCDEFGHIJKLMNOPQRSTUVWXYZ0123456789}

        \texttt{2BQF5WRTD8IJ6HLCOSUVK3A0X9YZN1G4ME7P}

        \begin{example} Under the above substitution cipher, what is the
        ciphertext for the message: ``March 12 0300''
        \end{example}
}{% presentation
\vspace{2in}
}{% nopresentation
\vspace{2in}
}{% comments
}

\content{
        \begin{example} What are the secret keys for a substitution cipher, and
        how many secret keys are there? 
        \end{example}
}{% presentation
\vspace{2in}
}{% nopresentation
\vspace{2in}
}{% comments
Unfortunately, this method is still not secure, due to frequency analysis (and
even frequency analysis with pairs of letters).
}

\subsection{Transposition Ciphers}

\content{
        \begin{example}
        Encrypt the message ``Meet at first and pine at midnight'' using a
        transposition cipher, and secret key ``help''.
        \end{example}
}{% presentation
\vspace{3in}
}{% nopresentation
\vspace{3in}
}{% comments
}

\content{
        \begin{example} The following is an encrypted message using a
        transposition cipher with secret key ``train''. What was the message?
        
        \texttt{RRSOGRANEI TZAAEOTTVN}
        \end{example}
}{% presentation
\vspace{3in}
}{% nopresentation
\vspace{3in}
}{% comments
}

\subsection{Advanced shared symmetric-key methods}


\content{
        Each method so far has been insecure due to frequency analysis. The next
        method gets around that vulnerability.

        \begin{example} Consider an encryption method using a shift cipher, but
        after each letter is encrypted, the shift is changed by some known
        amount. The secret key in this system is then the starting shift amount,
        and how the shift changes between each letter. Let's see an example,
        where the shift changed by +1 after each letter. Encrypt the phrase "See
        me". 

        \texttt{ABCDEFGHIJKLMNOPQRSTUVWXYZ0123456789}

        \texttt{2BQF5WRTD8IJ6HLCOSUVK3A0X9YZN1G4ME7P}
        \end{example}
}{% presentation
}{% nopresentation
}{% comments
}
